% Síntesis axiomática incorporada desde el anexo de caminos.
% Las pruebas operatorias extensas permanecen en sus fuentes.

\section{Reconstrucción de los principios cuánticos desde el estado genealógico}
\label{sec:postulados-cuanticos-genealogicos-rev2}
\index{mecánica cuántica!reconstrucción genealógica}
\index{estado!realización de Hilbert}

La formulación de Hilbert organiza la mecánica cuántica mediante estados,
observables autoadjuntos, evolución unitaria, composición tensorial y una
regla probabilística para los resultados~\cite{vonNeumann1932}. HMT--MD
construye una capa anterior a esa representación: una sección compatible,
su ruta observable, la memoria del transporte y el lector que publica un
resultado. El espacio de Hilbert no se reemplaza; aparece como codominio de
una realización lineal declarada.

La secuencia tipada es
\[
 \boxed{\begin{aligned}
 \text{sección genealógica}
 &\longrightarrow \text{realización lineal}
 \longrightarrow \text{fase y superposición}\\
 &\longrightarrow \text{observable y evolución}
 \longrightarrow \text{probabilidad y medición}
 \end{aligned}}
\]
Cada flecha conserva su dominio. En particular, el vector de estado no
sustituye a la sección TPK, del mismo modo que el valor real no sustituye a
la genealogía numérica que lo produce.

\subsection{Estado genealógico y realización de Hilbert}

\begin{definition}[Estado genealógico realizado]
Un estado genealógico es una quíntupla
\begin{equation}
 \mathfrak G=(x_\infty,\Gamma,\mathcal F,\mathcal M,\mathcal R),
 \label{eq:estado-genealogico-realizado-rev2}
\end{equation}
donde \(x_\infty\in\Omega_\infty\) es una sección compatible,
\(\Gamma\) su ruta observable, \(\mathcal F\) la familia de realizaciones,
\(\mathcal M\) la memoria de los levantamientos y \(\mathcal R\) el lector
declarado. Sea \(\mathsf P_{\rm TPK}\) la categoría cuyos objetos son las
secciones TPK preparadas, finitas o coinductivas, y cuyos morfismos son las
rutas registradas compatibles con sus datos de frontera. En particular, la
restricción de profundidad es un morfismo
\[
 r_{n+1,n}:x_{n+1}\longrightarrow x_n.
\]
Sea \(\mathsf{Hilb}_{\mathbb R,J}\) la categoría cuyos objetos
son pares \((\mathcal H_{\mathbb R},J)\), con \(\mathcal H_{\mathbb R}\)
espacio de Hilbert real y \(J\) ortogonal, \(J^2=-I\), y cuyos morfismos son
aplicaciones lineales reales acotadas \(T\) que entrelazan las estructuras
complejas, \(TJ_1=J_2T\). Una realización cuántica es un funtor
\[
 \mathcal Q:\mathsf P_{\rm TPK}\longrightarrow
 \mathsf{Hilb}_{\mathbb R,J},
\]
y una \emph{realización apuntada} es el par \((\mathcal Q,\Psi)\) que asigna
a cada sección preparada \(x\) un vector no nulo
\(\Psi_x\in\mathcal Q(x)\), compatible con las restricciones del sistema
inverso.
\end{definition}

Con la convención lineal en el segundo argumento, la forma hermítica asociada
se escribe
\[
 \langle u,v\rangle_{\mathbb C}
 =\langle u,v\rangle_{\mathbb R}
  -i\langle u,J_E v\rangle_{\mathbb R}.
\]
Elegir la carta \(J_E\leftrightarrow i\) realiza la estructura compleja que
en las Partes I--III se construye como raíz interna de \(-I\). Para prefijos de
profundidad \(n\), pongamos
\(\mathcal H_n:=\mathcal Q(x_n)\) y
\[
 p_{n+1,n}:=\mathcal Q(r_{n+1,n}):
 \mathcal H_{n+1}\longrightarrow\mathcal H_n.
\]
La compatibilidad entre restricción y realización exige
\begin{equation}
 p_{n+1,n}\Psi_{n+1}(x_{n+1})=\Psi_n(x_n).
 \label{eq:compatibilidad-realizacion-hilbert-rev2}
\end{equation}
Si \(r_{\infty,n}:x_\infty\to x_n\) es la restricción coinductiva y
\(p_n:=\mathcal Q(r_{\infty,n}):\mathcal H_\infty\to\mathcal H_n\), un
límite \(\mathcal H_\infty\) realiza la sección completa cuando
\(p_n\Psi_\infty(x_\infty)=\Psi_n(x_n)\) en cada profundidad.

\begin{principle}[Estado cuántico realizado]
El estado puro asociado a una sección se representa por el rayo del vector
no nulo \(\Psi=\Psi_\infty(x_\infty)\); una preparación estadística se
representa por un operador densidad positivo y de traza uno sobre la misma
realización. La identificación de fase global actúa en la representación:
la sección y su registro permanecen en el dominio genealógico.
\end{principle}

\subsection{Superposición y fase interna}

Si \(\Psi,\Phi\in\mathcal H_\infty\) y
\(a=\alpha+i\beta\in\mathbb C\), la acción de \(a\) en la presentación real
es \(\alpha I+\beta J_E\). La combinación \(a\Psi+b\Phi\) se forma después
de transportar las dos alternativas hasta una fibra común. Los levantamientos
genealógicos del \cref{chap:genealogia-operaciones-rev2} toman valores en
\(\text{salida}\times\text{memoria}\); pertenecen al dominio TPK y no se
multiplican como operadores. Si una realización proyectiva asigna operadores
unitarios \(U_g\) a transformaciones visibles, su defecto de composición se
describe por separado mediante
\[
 U_gU_h=c(g,h)U_{gh},
 \qquad c(g,h)\in U(1),
\]
con \(c(e,g)=c(g,e)=1\) y
\[
 c(g,h)c(gh,k)=c(h,k)c(g,hk).
\]
Un cociclo con valores en otro grupo de memoria sólo actúa en la realización
después de declarar una representación de ese grupo que entrelace los
operadores \(U_g\).

\begin{principle}[Superposición genealógica]
Las alternativas compatibles se transportan primero en
\(\mathsf P_{\rm TPK}\) y se combinan después linealmente en
\(\mathcal H_\infty\). La acción de \(J_E\) representa la fase relativa y
el peso de la ruta conserva el dato de elevación declarado.
\end{principle}

La interferencia se expresa mediante los términos cruzados
\[
 \|\Psi_1+\Psi_2\|^2
 =\|\Psi_1\|^2+\|\Psi_2\|^2
  +2\operatorname{Re}\langle\Psi_1,\Psi_2\rangle.
\]
Dos rutas con iguales extremos pueden, por tanto, contribuir con fases
distintas sin que su publicación conjunta borre la diferencia de historias.

\subsection{Observables y cálculo espectral}

Sea \(O:\Omega_\infty\to Y\subseteq\mathbb R\) un lector genealógico real.
Una realización cuántica de \(O\) es un operador autoadjunto
\(\widehat O\), con dominio denso \(D(\widehat O)\), cuyo PVM satisface la
condición de compatibilidad
\[
 P_O(B)\Psi_x
 \quad\hbox{representa la restricción de \(x\) a \(O^{-1}(B)\)}
\]
para cada boreliano \(B\subseteq\mathbb R\) en el soporte de la preparación.
En dimensión finita,
\[
 \widehat O=\sum_a o_a P_a,
 \qquad
 P_a P_b=\delta_{ab}P_a,
 \qquad
 \sum_a P_a=I;
\]
en dimensión general, la suma se sustituye por la integral espectral
\(\widehat O=\int\lambda\,\mathrm dP_O(\lambda)\).

\begin{principle}[Observable tipado]
El observable físico es el par formado por el lector en el dominio
genealógico y su representación autoadjunta. El cálculo espectral publica los
resultados posibles de la representación sin borrar la operación que los
definió.
\end{principle}

Para un estado puro normalizado \(\Psi\in D(\widehat O^2)\),
\[
 \langle O\rangle_\Psi
 =\langle\Psi,\widehat O\Psi\rangle,
 \qquad
 (\Delta_\Psi O)^2
 =\langle\Psi,(\widehat O-\langle O\rangle_\Psi)^2\Psi\rangle.
\]
La formulación mediante \(\operatorname{tr}(\rho\widehat O)\) extiende la
definición a estados mixtos cuando \(\rho\widehat O\) es de clase traza; el
segundo momento requiere asimismo que
\(\rho\widehat O^2\) sea de clase traza.

\subsection{Evolución, Hamiltoniano y ecuación de Schrödinger}

La evolución genealógica preserva la composición de transportes. Su
realización satisface
\[
 U(t_3,t_2)U(t_2,t_1)=U(t_3,t_1),
 \qquad U(t,t)=I.
\]
Cuando preserva el producto interno y satisface
\(U(t)J_E=J_EU(t)\), es ortogonal en la carta real y unitario en la carta
compleja. Para un grupo unitario uniparamétrico fuertemente continuo, el
teorema de Stone proporciona un generador autoadjunto y \(J_E\)-lineal
\(H\)~\cite{Stone1932}:
\[
 U(t)=\exp\!\left(-J_E H t/\hbar_{\rm int}\right).
\]
Sobre el dominio \(D(H)\), la derivación temporal da
\begin{equation}
 \hbar_{\rm int}J_E\dot\Psi(t)=H\Psi(t),
 \label{eq:schrodinger-genealogico-rev2}
\end{equation}
y la carta \(J_E\leftrightarrow i\) produce
\(i\hbar_{\rm int}\dot\Psi=H\Psi\).

El calendario \(C_{108}\) realiza esta correspondencia de manera finita:
el Hamiltoniano autoadjunto construido en su fuente tiene por
exponencial, durante \(t_0\), el desplazamiento unitario de una iteración. El
grupo continuo requiere adicionalmente continuidad fuerte; no se infiere
sólo de la periodicidad del calendario.

\subsection{Probabilidad y representación cuadrática de las ramas}

Sea \(\mathcal C_n=\{C_s:s\in\mathcal S_n\}\) la partición de
\(\Omega_\infty\) por un conjunto finito \(\mathcal S_n\) de prefijos de
profundidad \(n\), y sea \(\mu\) una medida cilíndrica de probabilidad
compatible, \(\mu(\Omega_\infty)=1\). Sea \(E\) un espacio de Hilbert real
con estructura compleja \(J_E\). En
\(\mathcal H_n=\ell^2(\mathcal S_n)\otimes E\), se define
\begin{equation}
 \Psi_n=\sum_{s\in\mathcal S_n}
 \sqrt{\mu(C_s)}\,|s\rangle\otimes e^{\theta_s J_E}u_0,
 \label{eq:estado-amplitudes-medida-rev2}
\end{equation}
con \(u_0\) unitario.

\begin{theorem}[Representación cuadrática de una medida de ramas]
\label{thm:born-cilindrico-rev2}
Para un suceso cilíndrico
\(A=\bigcup_{s\in I_A}C_s\), sea
\[
 P_A=\left(\sum_{s\in I_A}|s\rangle\langle s|\right)\otimes I_E.
\]
Entonces
\[
 \|P_A\Psi_n\|^2=\mu(A).
\]
\end{theorem}

\begin{proof}
Los vectores de prefijo son ortogonales y
\(e^{\theta_s J_E}\) preserva la norma. Por tanto,
\[
 \|P_A\Psi_n\|^2
 =\sum_{s\in I_A}\mu(C_s)=\mu(A).
\]
\end{proof}

La regla cuadrática se obtiene aquí para el PVM de sucesos cilíndricos de
profundidad finita como representación exacta de la medida de ramas
declarada~\cite{Born1926}. En el sentido inverso, el teorema de Gleason
caracteriza, en dimensión al menos tres, las medidas aditivas no contextuales
sobre proyectores mediante operadores densidad~\cite{Gleason1957}. El
acoplamiento con el aparato y el condicionamiento del resultado se desarrollan
después en la teoría de la medición.

\subsection{Incompatibilidad de lectores e incertidumbre}

Las traslaciones y fases nonádicas definen el par de Weyl
\[
 \mathcal H_9=\ell^2(\mathbb Z/9\mathbb Z),
 \qquad
 \omega_9=e^{2\pi J_E/9},
 \qquad
 X|n\rangle=|n+1\rangle,
 \qquad
 Z|n\rangle=\omega_9^n|n\rangle.
\]
Los índices se leen módulo \(9\), y \(\omega_9\) es el escalar complejo
realizado por la estructura \(J_E\) sobre \(\mathcal H_9\).
La identidad
\[
 ZX=\omega_9XZ
\]
realiza el grupo finito de Heisenberg de orden \(9^3=729\). Este cardinal
pertenece al grupo de fases y traslaciones y conserva su tipo frente a las
demás apariciones de \(729\).

Para observables autoadjuntos \(A\) y \(B\), un vector normalizado
\(\Psi\in D(AB)\cap D(BA)\) con segundos momentos finitos permite aplicar
Cauchy--Schwarz y obtener la
desigualdad de Robertson--Heisenberg
\begin{equation}
 \Delta_\Psi A\,\Delta_\Psi B
 \geq\frac12
 \bigl|\langle\Psi,[A,B]\Psi\rangle\bigr|.
 \label{eq:robertson-hmt-rev2}
\end{equation}
La incompatibilidad se presenta así en niveles coordinados: el conmutador de
las evaluaciones de APP, la relación de Weyl en la representación y la cota
de dispersión para las lecturas físicas~\cite{Heisenberg1927,Robertson1929}.

Para \(\Psi\neq0\), la transformada discreta de Fourier aporta la cota de soporte
\[
 |\operatorname{supp}\Psi|\,
 |\operatorname{supp}\widehat\Psi|\geq9,
\]
y, para bases mutuamente insesgadas de posición y fase, la desigualdad
entrópica de Shannon, con logaritmos naturales,
\[
 H_Q(\Psi)+H_P(\Psi)\geq\log9
\]
ofrece una tercera formulación exacta~\cite{MaassenUffink1988}. Aquí
\(H_Q\) y \(H_P\) son las entropías de Shannon de las distribuciones de
probabilidad en las dos bases.

\subsection{Composición y entrelazamiento}

La composición genealógica pega secciones sobre datos de frontera. Fijemos
en \(\mathsf P_{\rm TPK}\) una estructura monoidal
\((\boxtimes,\mathbf 1,a,l,r)\), con producto, unidad y coherencias de
asociatividad y unidad declaradas. Una realización monoidal fuerte
proporciona isomorfismos naturales
\[
 \mu_{A,B}:
 \mathcal Q(A)\widehat\otimes\mathcal Q(B)
 \xrightarrow{\sim}\mathcal Q(A\boxtimes B).
\]
En la imagen de Hilbert aparece
\(\mathcal H_A\otimes\mathcal H_B\). Un estado puro no separable admite una
descomposición de Schmidt
\[
 |\Psi\rangle
 =\sum_r\sqrt{p_r}\,|r\rangle_A\otimes|r\rangle_B,
\]
y cumple \(p_r\geq0\), \(\sum_rp_r=1\), con
\(\{|r\rangle_A\}\) y \(\{|r\rangle_B\}\) bases ortonormales de Schmidt,
y sus reducciones poseen el mismo espectro no nulo. La entropía asociada a
ese espectro recibe, en el capítulo de Barbero--Immirzi, una realización de
área y borde.

\begin{principle}[Composición genealógica]
El producto tensorial representa un pegado de secciones sobre una frontera
compartida. El entrelazamiento expresa la no factorización de la sección
realizada respecto de la bipartición elegida.
\end{principle}

\subsection{Teorema de reconstrucción cinemática}

\begin{theorem}[Reconstrucción hilbertiana condicionada del estado genealógico]
\label{thm:cinematica-cuantica-tpk-rev2}
Sea \(\Omega_\infty\) un espacio de secciones compatibles, sea
\((\mathsf P_{\rm TPK},\boxtimes,\mathbf 1)\) la categoría monoidal de rutas
con las coherencias declaradas y sea
\[
 \mathcal Q:\mathsf P_{\rm TPK}\longrightarrow
 \mathsf{Hilb}_{\mathbb R,J}
\]
una realización monoidal fuerte, apuntada por una sección compatible
\(\Psi_x\in\mathcal Q(x)\). Supóngase que:
\begin{enumerate}
 \item \(J_E^2=-I\), \(J_E\) preserva el producto interno y los morfismos de
 \(\mathcal Q\) entrelazan las estructuras complejas;
 \item cada lector real declarado se representa por un operador autoadjunto
 cuyo PVM satisface la condición de compatibilidad con las fibras del lector;
 \item cada evolución continua se representa por un grupo unitario
 uniparamétrico fuertemente continuo \(U(t)\), con
 \(U(t)J_E=J_EU(t)\), y generador autoadjunto \(J_E\)-lineal \(H\);
 \item la preparación induce una medida cilíndrica de probabilidad compatible
 y el estado \(\Psi_n\) de \eqref{eq:estado-amplitudes-medida-rev2}.
\end{enumerate}
Entonces la carta compleja de \(\mathcal Q\) realiza estados en rayos o
densidades, superposición compleja, observables autoadjuntos, evolución
unitaria generada por \(H\), composición tensorial y probabilidades
cuadráticas para el PVM cilíndrico de profundidad finita.
\end{theorem}

\begin{proof}
La primera hipótesis convierte cada fibra real en un espacio de Hilbert
complejo y hace complejamente lineales los transportes. La segunda permite
aplicar el teorema espectral con el lector tipado. La tercera proporciona el
Hamiltoniano por el teorema de Stone y la ecuación diferencial sobre
\(D(H)\). La cuarta y el \cref{thm:born-cilindrico-rev2} dan la regla
cuadrática sobre los sucesos cilíndricos. La monoidalidad fuerte realiza el
producto tensorial y su noción de separabilidad.
\end{proof}

\begin{statusbox}[title={Resultado y alcance}]
La reconstrucción hilbertiana queda coordinada como propiedad de una
realización del estado genealógico: la estructura compleja procede de
\(J_E^2=-I\); la evolución, de la representación de rutas; la probabilidad,
de una medida compatible; y el entrelazamiento, de la composición sobre una
frontera compartida. El teorema es exacto bajo sus cuatro hipótesis y no
identifica automáticamente TPK con una teoría cuántica concreta. La
construcción de una realización física concreta conserva sus fuentes:
Hamiltoniano, aparato, lectores, dominio y condiciones de borde.
\end{statusbox}
